% SHARD-ID: AQM-07-TORIC-OPERATOR-RELATION
% SHARD-TITLE: Operator Relation Between Cellular and Ghost Supercharges
% SHARD-SUMMARY: Shows that the cellular and ghost supercharges are not restrictions of one operator.
% SHARD-SUMMARY: Derives the ghost projectors from the boundary maps via translations and syndrome characters.
% SHARD-SUMMARY: Proves that degree-zero ghost cohomology is the complex linearization of cellular cohomology.
% SHARD-KEYWORDS: toric code, supercharge, operator relation, ghost Fock space, cellular differential, Koszul

\section{Operator Relation Between Cellular and Ghost Supercharges}
\label{sec:operator-relation}

The honest operator-level answer is: \(Q_{\mathrm{cell}}\) is not the
restriction of \(Q_{\mathrm{gh}}\) to a low-ghost-number subspace of the
check-ghost Fock space. The precise relation is induced, not restrictive:
the cellular maps \(\delta^0,\delta^1\) define translations and syndrome
characters on the physical edge-bit Hilbert space, and \(Q_{\mathrm{gh}}\) is
the Koszul differential for the corresponding commuting projector equations.

\subsection{Why It Is Not a Low-Ghost Restriction}

Let
\[
  E=C^1
\]
be the \(\mathbb F_2\)-vector space of edge-bit strings, and let
\[
  \mathcal H=\mathbb C[E]
\]
be the complex vector space with basis \(\{|a\rangle:a\in E\}\). The
check-ghost space used by \(Q_{\mathrm{gh}}\) is
\[
  G_{\mathrm{check}}
  =\mathbb C\{\eta_v:v\in C_0\}
  \oplus \mathbb C\{\theta_f:f\in C_2\}.
\]
The full ghost-extended space is
\[
  \mathcal H_{\mathrm{gh}}
  =
  \mathcal H\otimes\Lambda(G_{\mathrm{check}}).
\]
Its low ghost-number pieces are therefore
\[
  \mathcal H\otimes\Lambda^0G_{\mathrm{check}}=\mathcal H,
\]
\[
  \mathcal H\otimes\Lambda^1G_{\mathrm{check}}
  =
  \mathcal H\otimes
  \bigl(\mathbb C\{v\text{-ghosts}\}\oplus
        \mathbb C\{f\text{-ghosts}\}\bigr).
\]
There is no edge one-particle sector here. But the cellular differential has
the form
\[
  Q_{\mathrm{cell}}:\ C^0\to C^1\to C^2.
\]
Its middle degree is the edge space \(C^1\). In contrast,
\[
  Q_{\mathrm{gh}}:\ \mathcal H\otimes\Lambda^rG_{\mathrm{check}}
  \longrightarrow
  \mathcal H\otimes\Lambda^{r+1}G_{\mathrm{check}}
\]
adds a vertex-check or face-check ghost and keeps the same physical Hilbert
space. On ghost degree zero it maps
\[
  \mathcal H\longrightarrow
  \mathcal H\otimes(\mathbb C\{v\text{-ghosts}\}\oplus
                    \mathbb C\{f\text{-ghosts}\}),
\]
not \(C^0\to C^1\). On ghost degree one it maps to two-check ghosts, not to
faces. Thus \(Q_{\mathrm{cell}}\) is not literally \(Q_{\mathrm{gh}}\) on any
canonical low-ghost subspace of \(\mathcal H_{\mathrm{gh}}\).

If one wants a cell-fermion operator whose one-particle restriction is
\(Q_{\mathrm{cell}}\), one must build a different Fock space. Choose an
oriented signed cellular boundary over \(\mathbb Z\),
\[
  C_2^{\mathbb Z}\xrightarrow{\partial_2^{\mathbb Z}}
  C_1^{\mathbb Z}\xrightarrow{\partial_1^{\mathbb Z}} C_0^{\mathbb Z},
  \qquad
  \partial_1^{\mathbb Z}\partial_2^{\mathbb Z}=0.
\]
Let \(W=W_0\oplus W_1\oplus W_2\) be the complex vector space with basis the
oriented cells, and let \(a_\sigma^*,a_\sigma\) be creation and annihilation
operators for cell fermions. The number-preserving CAR bilinear with this
one-particle matrix is
\[
  \widehat Q_{\mathrm{cell}}
  =
  \sum_{e,v}(\delta^0_{\mathbb Z})_{ev}a_e^*a_v
  +
  \sum_{f,e}(\delta^1_{\mathbb Z})_{fe}a_f^*a_e .
\]
It preserves fermion number and raises the cell degree of each one-particle
label. On the one-particle sector \(W\), its restriction is exactly the signed
cellular cochain differential and squares to zero because
\(\delta^1_{\mathbb Z}\delta^0_{\mathbb Z}=0\). No full-Fock nilpotence claim is
made for this number-preserving bilinear. Such a claim would require a
separate graded exterior-chain derivation with Koszul signs, which is not fixed
in the local conventions here. This is therefore only a one-particle
cell-fermion comparison, not a check-ghost construction and not the
check-ghost Fock space used by \(Q_{\mathrm{gh}}\). Reducing the signed
incidence matrices modulo \(2\) gives the CSS supports used below.

\subsection{The Induced Operator Construction}

The actual bridge starts with the mod-\(2\) cellular maps
\[
  C^0\xrightarrow{\delta^0}E\xrightarrow{\delta^1}C^2,
  \qquad \delta^1\delta^0=0.
\]
The edge-bit space \(E\) is an abelian group under addition. For each
\(u\in E\), define the translation operator on \(\mathcal H=\mathbb C[E]\) by
\[
  T_u|a\rangle=|a+u\rangle .
\]
For each \(s\in E\), define the character multiplication operator by
\[
  M_s|a\rangle=(-1)^{s\cdot a}|a\rangle .
\]
These obey
\[
  T_uM_s=(-1)^{s\cdot u}M_sT_u.
\]

Now feed the boundary maps into these two representations of \(E\). For a
vertex \(v\), set
\[
  A_v=T_{\delta^0 v}.
\]
For a face \(f\), set \(s_f=\partial_2 f\in E\) and
\[
  B_f=M_{s_f}.
\]
Then
\[
  A_vB_f=(-1)^{s_f\cdot\delta^0v}B_fA_v.
\]
The exponent \(s_f\cdot\delta^0v\) is the \(f\)-coordinate of
\(\delta^1\delta^0v\), equivalently the \(v\)-coefficient of
\(\partial_1\partial_2f\). It is zero. Hence all star and plaquette checks
commute.

The violated-check projectors are therefore
\[
  P_v^X=\frac{I-T_{\delta^0v}}{2},\qquad
  P_f^Z=\frac{I-M_{s_f}}{2}.
\]
Their action on the computational basis is concrete:
\[
  P_v^X|a\rangle
  =
  \frac{|a\rangle-|a+\delta^0v\rangle}{2},
\]
and
\[
  P_f^Z|a\rangle
  =
  \frac{1-(-1)^{s_f\cdot a}}{2}|a\rangle
  =
  (\delta^1a)_f|a\rangle,
\]
where the bit \((\delta^1a)_f\in\mathbb F_2\) is embedded as \(0\) or \(1\) in
\(\mathbb C\). Thus the plaquette part of \(Q_{\mathrm{gh}}\) reads off the
cellular coboundary \(\delta^1a\) as a face-ghost syndrome, while the star part
is a finite-difference operator along the \(\operatorname{im}\delta^0\) orbit.

\subsection{The Ghost Differential}

Let exterior multiplication by \(\eta_v\) and \(\theta_f\) be denoted by the
same symbols. Define
\[
  Q_X=\sum_v \eta_v P_v^X,\qquad
  Q_Z=\sum_f \theta_f P_f^Z,\qquad
  Q_{\mathrm{gh}}=Q_X+Q_Z .
\]
The projectors commute because the checks commute, and exterior
multiplication by ghost labels anticommutes. Therefore
\[
  Q_{\mathrm{gh}}^2
  =
  \sum_{i,j}\gamma_i\gamma_jP_iP_j
  =
  \sum_{i<j}(\gamma_i\gamma_j+\gamma_j\gamma_i)P_iP_j
  =0,
\]
where \(\gamma_i\) ranges over all \(\eta_v,\theta_f\). With ghost
annihilation operators \(\gamma_i^*\) satisfying
\(\gamma_i^*\gamma_j+\gamma_j\gamma_i=\delta_{ij}\), one also has
\[
  Q_{\mathrm{gh}}Q_{\mathrm{gh}}^*
  +Q_{\mathrm{gh}}^*Q_{\mathrm{gh}}
  =\sum_iP_i .
\]
Thus \(Q_{\mathrm{gh}}\) is the Koszul differential for the equations
\[
  T_{\delta^0v}\psi=\psi,\qquad M_{s_f}\psi=\psi .
\]
Those equations are exactly the stabilizer equations.

\subsection{Cohomology-Level Relation}

Let
\[
  \psi=\sum_{a\in E}\psi_a|a\rangle\in\mathcal H
\]
be a no-ghost state. The plaquette equations are:
\[
  P_f^Z\psi=0\ \text{for all }f
  \quad\Longleftrightarrow\quad
  \psi_a=0\ \text{unless}\ \delta^1a=0.
\]
Indeed, \(P_f^Z\) multiplies the coefficient of \(|a\rangle\) by the bit
\((\delta^1a)_f\). The star equations are:
\[
  P_v^X\psi=0\ \text{for all }v
  \quad\Longleftrightarrow\quad
  \psi_a=\psi_{a+\delta^0v}\ \text{for all }a,v.
\]
Thus a no-ghost state is \(Q_{\mathrm{gh}}\)-closed exactly when its
coefficients are supported on \(\ker\delta^1\) and are constant on cosets of
\(\operatorname{im}\delta^0\). Since there are no degree \(-1\) states, there
are no no-ghost exact states. Therefore
\[
  H^0(Q_{\mathrm{gh}})
  \cong
  \mathbb C[\ker\delta^1/\operatorname{im}\delta^0]
  =
  \mathbb C[H^1(Q_{\mathrm{cell}})].
\]

This is the concrete relation. \(Q_{\mathrm{cell}}\) does not act inside the
check-ghost Fock space. Instead, its two arrows induce two kinds of operators
on \(\mathbb C[C^1]\): \(\delta^0\) gives translations, and \(\delta^1\) gives
syndrome multiplications. The ghost supercharge is the Koszul operator for
the projector constraints built from those induced operators. The equality is
therefore a cohomology-level identification, not an operator restriction:
\[
  H^0(Q_{\mathrm{gh}})
  \simeq
  \mathbb C[H^1(Q_{\mathrm{cell}})].
\]
